% Part: model-theory
% Chapter: models-of-arithmetic
% Section: introduction

\documentclass[../../../include/open-logic-section]{subfiles}

\begin{document}

\olfileid{mod}{mar}{int}
\section{مقدمة}

\emph{النموذج القياسي} للحساب هو الـ!!{structure}~$\Struct{N}$ ذات
$\Domain{N} = \Nat$، التي تفسر فيها $\Obj{0}$ و$\prime$ و$+$ و$\times$
و$<$ على النحو المتوقع. أي إن $\Obj{0}$ هو~$0$، و$\prime$ دالة الخلف،
ويفسر $+$ بجمع الأعداد في~$\Nat$، و$\times$ بضربها، و$<$ بترتيبها
المعتاد. وعلى وجه التحديد:
\begin{align*}
  \Assign{\Obj{0}}{N} & = 0\\
  \Assign{\prime}{N}(n) & = n + 1\\
  \Assign{+}{N}(n, m) & = n + m\\
  \Assign{\times}{N}(n, m) & = nm\\
  \Assign{<}{N} & = \Setabs{\tuple{n,m} \in \Nat^2}{n<m}
\end{align*}
وبالطبع توجد بنى للغة $\Lang{L_A}$ مجالاتها غير~$\Nat$. فمثلًا يمكن أن
نأخذ $\Struct{M}$ ذات المجال $\Domain{M} = \{a\}^*$ (أي المتتاليات
المنتهية للرمز الواحد~$a$، وهي $\emptyset$، $a$، $aa$، $aaa$، \dots)،
والتفسيرات:
\begin{align*}
  \Assign{\Obj{0}}{M} & = \emptyset\\
  \Assign{\prime}{M}(s) & = s \concat a\\
  \Assign{+}{M}(s, t) & = s \concat t\\
  \Assign{\times}{M}(s, t) & = a^{\lvert s\rvert\lvert t\rvert}\\
  \Assign{<}{M} & = \Setabs{\tuple{s,t} \in (\{a\}^*)^2}{\lvert s\rvert<\lvert t\rvert}
\end{align*}
هاتان البنيتان «الشيء نفسه في جوهرهما»، بمعنى أن الفرق الوحيد بينهما في
!!{element}s الـ!!{domain}s، لا في الكيفية التي تربط بها دوال التفسير
!!{element}s الـ!!{domain}s بعضها ببعض. ونقول إن الـ!!{structure}s
\emph{متشاكلة}.

من النتائج السهلة لمبرهنة التراص أن لكل نظرية صادقة في~$\Struct{N}$
نماذج غير متشاكلة مع~$\Struct{N}$. وتسمى هذه البنى \emph{غير قياسية}.
واللافت فيها أن !!{element}s النموذج القياسي (أي $\Struct{N}$، وكذلك كل
!!{structure}s المتشاكلة معها) تستنفدها قيم الأعداد القياسية~$\num{n}$،
أي:
\[
\Domain{N} = \Setabs{\Value{\num{n}}{N}}{n \in \Nat}
\]
أما في النماذج غير القياسية فليس الأمر كذلك: إذا كانت $\Struct{M}$ غير
قياسية، وجد على الأقل $x \in \Domain{M}$ بحيث
$x \neq \Value{\num{n}}{M}$ لكل~$n$.

وهذه العناصر غير القياسية طريفة حقًا: فهي «أعداد طبيعية لا نهائية».
كما أن وجودها يفسر، بمعنى ما، ظواهر عدم الاكتمال. خذ مثلًا عبارة اتساق
حساب بيانو، $\OCon[\Th{PA}]$، أي
$\lnot \lexists[x][\OPrf[\Th{PA}](x, \gn{\lfalse})]$. ولما كانت
$\Th{PA}$ لا تثبت $\OCon[\Th{PA}]$ ولا $\lnot \OCon[\Th{PA}]$، أمكن
إضافة أي منهما إلى $\Th{PA}$ مع بقاء الاتساق. وبما أن $\Th{PA}$ متسقة،
فإن $\Sat{N}{\OCon[\Th{PA}]}$، وبالتالي
$\Sat/{N}{\lnot \OCon[\Th{PA}]}$. إذن \emph{ليست} $\Struct{N}$ نموذجًا لـ
$\Th{PA} \cup \{\lnot \OCon[\Th{PA}]\}$، ولا بد أن تكون جميع نماذج
هذه النظرية غير قياسية. ويجب أن تحتوي نماذج
$\Th{PA} \cup \{\lnot \OCon[\Th{PA}]\}$ بعض الـ!!{element} الذي يؤدي
دور الشاهد الذي يجعل
$\lexists[x][\OPrf[\Th{PA}](x,\gn{\lfalse})]$ صادقة، أي عدد غودل
لـ!!a{derivation} لتناقض من~$\Th{PA}$. ولا يمكن أن يكون هذا
الـ!!a{element} قياسيًا، لأن
$\Th{PA} \Proves \lnot \OPrf[\Th{PA}](\num{n}, \gn{\lfalse})$
لكل~$n$.

\end{document}
